% ============================================================
%  Errades del solucionari — Matemàtiques II (2n de Batxillerat)
% ============================================================

% ---------- OPCIONS ----------
% De cada parella, deixa sense % només UNA de les dues línies.

% Es veu el «Motiu» de cada errada?
\newif\ifmotiu \motiufalse     % NO: només «On diu» i «Ha de dir»
%\newif\ifmotiu \motiutrue     % SÍ: també el «Motiu»

% Es veuen les Observacions (comentaris que no són errades)?
\newif\ifobs \obsfalse     % NO
%\newif\ifobs \obstrue     % SÍ

% ------------------------------------------------------------
%  main.tex és el fitxer principal: només crida els altres.
%    headers.tex            paquets i format de pàgina
%    defs.tex               comandes i entorns propis
%    errors-setmanaN.tex    les errades de la setmana N
%
%  PER AFEGIR UNA SETMANA
%    1. Copia errors-setmana4.tex amb el nom errors-setmana5.tex
%       i canvia'n el contingut.
%    2. Afegeix al final d'aquest fitxer, abans de \end{document}:
%           \newpage
%           % ============================================================
%  Setmana 5 — Unitat 8 · Derivades
%  Exercicis: 37, 86, 87, 88, 92
% ============================================================

\setmana{5}{Unitat 8 \textperiodcentered{} Derivades}{18 d'octubre de 2026}%
        {37, 86, 87, 88 i 92}

\noindent
No hi ha errors.

%\senseerrades{37, 86, 87, 88 i 92}

% ---------- Observacions (només surten amb \obstrue) ----------
\ifobs
\bigskip
\noindent
Una observació, que no és cap errada:

\begin{observacio}{Exercici 87}
  L'enunciat demana calcular la funció derivada «utilitzant la definició»,
  però el solucionari només hi dona el resultat final de cada apartat, sense
  el desenvolupament. Els resultats són correctes. Si vols veure un model del
  desenvolupament pas a pas, mira la solució de l'exercici 88.
\end{observacio}
\fi

%
%  PER FER EL PDF D'UNA SOLA SETMANA
%    Posa un % davant dels \input (i dels \newpage) de les altres.
% ------------------------------------------------------------

\documentclass[11pt,a4paper]{article}

% ============================================================
%  headers.tex — paquets i format de pàgina
%  (el mateix estil que el full «Exemples de funcions no
%   contínues en un punt»)
% ============================================================

% ---------- Paquets ----------
\usepackage[utf8]{inputenc}
\usepackage[T1]{fontenc}
\usepackage[catalan]{babel}
\usepackage{amsmath,amssymb,amsthm}
\usepackage[margin=2.5cm,headheight=14pt]{geometry}
\usepackage{tcolorbox}
\usepackage{enumitem}
\usepackage{fancyhdr}
\usepackage[hidelinks]{hyperref}

% ---------- Capçalera i peu de pàgina ----------
% \titoldoc i \curs es defineixen a defs.tex
\pagestyle{fancy}
\fancyhf{}
\fancyhead[L]{\small\titoldoc}
\fancyhead[R]{\small\curs}
\fancyfoot[C]{\thepage}
\renewcommand{\headrulewidth}{0.4pt}

% ============================================================
%  defs.tex — comandes i entorns propis
% ============================================================

% ---------- Textos generals ----------
\newcommand{\titoldoc}{Errades del solucionari}
\newcommand{\curs}{Matemàtiques II \textperiodcentered{} 2n de Batxillerat}

% ---------- Notació ----------
\newcommand{\R}{\mathbb{R}}
\newcommand{\dlim}{\displaystyle\lim}   % límit amb el subíndex a sota

% ---------- Colors ----------
\colorlet{fonscaixa}{blue!4}
\colorlet{marccaixa}{gray!100!black}
\colorlet{colorondiu}{red!65!black}
\colorlet{colorhadedir}{green!40!black}

% ---------- Capçalera de cada setmana ----------
%   \setmana{número}{unitat}{data límit}{exercicis de la setmana}
\newcommand{\setmana}[4]{%
  \section*{Setmana #1}%
  \addcontentsline{toc}{section}{Setmana #1}%
  \noindent\textbf{#2}\quad\textbullet\quad data límit: #3\par
  \noindent Exercicis de la setmana: #4.\par
  \medskip}

% ---------- Una errada ----------
%   \begin{errada}{Exercici 94 b)}
%     \ondiu{el que hi ha imprès al solucionari}
%     \hadedir{el que hi hauria de dir}
%     \motiu{explicació breu}
%   \end{errada}
% El «Motiu» només surt al PDF si la flag de dalt de main.tex és \motiutrue;
% amb \motiufalse, el text de \motiu{...} s'ignora.
\newenvironment{errada}[1]{%
  \begin{tcolorbox}[colback=fonscaixa,colframe=marccaixa,
    title={#1},fonttitle=\bfseries]
  \begin{description}[style=multiline,leftmargin=2.3cm,
    font=\normalfont\bfseries,itemsep=3pt,topsep=0pt,parsep=0pt]}%
 {\end{description}\end{tcolorbox}}

\newcommand{\ondiu}[1]{\item[{\color{colorondiu}On diu}] #1}
\newcommand{\hadedir}[1]{\item[{\color{colorhadedir}Ha de dir}] #1}
\newcommand{\motiu}[1]{\ifmotiu\item[{\color{marccaixa}Motiu}] #1\fi}

% ---------- Una observació (no és cap errada, però pot confondre) ----------
%   \begin{observacio}{Exercici 120} ... \end{observacio}
% Al fitxer de cada setmana, el bloc d'observacions va entre \ifobs i \fi:
% només surt al PDF si la flag de dalt de main.tex és \obstrue.
\newtcolorbox{observacio}[1]{colback=gray!6,colframe=gray!60!black,
  title={Observació \textperiodcentered{} #1},fonttitle=\bfseries}

% ---------- Exercicis revisats i sense errades ----------
%   \senseerrades{103, 40 i 106}
% Als fitxers de les setmanes aquesta línia hi és escrita però comentada (%):
% per fer-la sortir al PDF, treu-li el %.
\newcommand{\senseerrades}[1]{%
  \medskip\noindent$\checkmark$\ \ \textbf{Revisats i sense errades:} #1.\par}

\hypersetup{pdftitle={\titoldoc}}

\title{\titoldoc\\[2mm]\large\curs}
\author{}
\date{}

\begin{document}
\maketitle
\vspace{-1.5cm}

% ---------- Setmanes ----------
% ============================================================
%  Setmana 1 — Unitat 7 · Límits i continuïtat
%  Exercicis: 44, 45, 46, 47, 48
% ============================================================

\setmana{1}{Unitat 7 \textperiodcentered{} Límits i continuïtat}{20 de setembre de 2026}%
        {44, 45, 46, 47 i 48}

\noindent
No hi ha errors.

%\senseerrades{44, 45, 46, 47 i 48}

% ============================================================
%  Setmana 2 — Unitat 7 · Límits i continuïtat
%  Exercicis: 66, 68, 70, 76, 88, 90
% ============================================================

\setmana{2}{Unitat 7 \textperiodcentered{} Límits i continuïtat}{27 de setembre de 2026}%
        {66, 68, 70, 76, 88 i 90}

\begin{errada}{Exercici 76 f)}
  \ondiu{$=\dlim_{x\to 3}\dfrac{(3-x)(3+x)}{(x-3)(x+2)}
    =\dlim_{x\to 3}\dfrac{3-x}{x+2}=0$}
  \hadedir{$=\dlim_{x\to 3}\dfrac{(3-x)(3+x)}{(x-3)(x+2)}
    =\dlim_{x\to 3}\dfrac{-(3+x)}{x+2}=-\dfrac{6}{5}$}
  \motiu{Els factors que se simplifiquen són $3-x$ i $x-3$, que són oposats:
    $\dfrac{3-x}{x-3}=-1$. Queda $-\dfrac{3+x}{x+2}$, que en $x=3$ val
    $-\dfrac{6}{5}$.}
\end{errada}

%\senseerrades{66, 68, 70, 88 i 90}

% ---------- Observacions (només surten amb \obstrue) ----------
\ifobs
\bigskip
\noindent
Una observació menor:\par\nopagebreak

\begin{observacio}{Exercici 68}
  A la solució, els apartats c) i d) surten en ordre canviat: primer el d) i
  després el c). Els valors de cadascun són correctes.
\end{observacio}
\fi

% ============================================================
%  Setmana 3 — Unitat 7 · Límits i continuïtat
%  Exercicis: 92, 93, 94, 102, 103, 40, 106
% ============================================================

\setmana{3}{Unitat 7 \textperiodcentered{} Límits i continuïtat}{4 d'octubre de 2026}%
        {92, 93, 94, 102, 103, 40 i 106}

\begin{errada}{Exercici 92 b)}
  \ondiu{$\dlim_{x\to 0^-} f(x)=\dlim_{x\to 0^+} f(x)=0{,}5$ \quad i \quad $f(2)=2{,}5$}
  \hadedir{$\dlim_{x\to 0^-} f(x)=\dlim_{x\to 0^+} f(x)=1$ \quad i \quad $f(2)=5$}
  \motiu{A la gràfica de l'enunciat cada quadre és una unitat (a l'eix $Y$,
    el $2$ és a la segona marca): el punt buit és a alçada $1$ i la corba
    passa pel punt $(2,5)$. És la mateixa paràbola de l'apartat d), on el
    solucionari sí que llegeix $f(-1)=2$ i límit $5$ en $x=2$.
    La conclusió no canvia: discontinuïtat evitable en $x=0$ i funció
    contínua en $x=2$.}
\end{errada}

\begin{errada}{Exercici 92 e)}
  \ondiu{«La funció no és contínua en $x=-2$; té una discontinuïtat de salt finit.»}
  \hadedir{«La funció no és contínua en $x=2$; té una discontinuïtat de salt finit.»}
  \motiu{És la frase que tanca l'estudi del punt $x=2$, on $f(2)=3$ i els
    límits laterals valen $-1$ i $3$. El punt $x=-2$ ja s'ha estudiat
    abans, i hi ha un salt infinit.}
\end{errada}

\begin{errada}{Exercici 93 c)}
  \ondiu{contínua en el seu domini, $(-\infty,-2)\cup(2,+\infty)$}
  \hadedir{contínua en el seu domini, $(-\infty,-2]\cup[2,+\infty)$}
  \motiu{La condició és $x^2-4\geq 0$, que també es compleix en $x=-2$ i en
    $x=2$: allà la funció existeix i val $\sqrt{0}=0$. La mateixa solució
    ho escriu bé just abans ($x\geq 2$, $x\leq -2$).}
\end{errada}

\begin{errada}{Exercici 93 g)}
  \ondiu{$\dlim_{x\to 1} 2\,|x-1| = 0 = y(0)$}
  \hadedir{$\dlim_{x\to 1} 2\,|x-1| = 0 = y(1)$}
  \motiu{El punt que s'estudia és $x=1$, i $y(1)=2\,|1-1|=0$.
    En canvi, $y(0)=2$.}
\end{errada}

\begin{errada}{Exercici 94 b)}
  \ondiu{$x^2-2x=3\neq 0$ per a qualsevol $x$ real}
  \hadedir{$x^2-2x+3\neq 0$ per a qualsevol $x$ real}
  \motiu{El denominador és $x^2-2x+3$, i no s'anul·la mai perquè el
    discriminant és negatiu: $(-2)^2-4\cdot 1\cdot 3=-8<0$.
    Per això la funció és contínua en tot $\R$.}
\end{errada}

\begin{errada}{Exercici 94 e), punt $x=-3$}
  \ondiu{«No existeix $\dlim_{x\to 3} f(x)$.»}
  \hadedir{«No existeix $\dlim_{x\to -3} f(x)$.»}
  \motiu{Els dos límits laterals que s'acaben de calcular són els del punt $x=-3$.}
\end{errada}

\begin{errada}{Exercici 94 e), punt $x=1$}
  \ondiu{«La funció té una discontinuïtat evitable en $x=3$.»}
  \hadedir{«La funció té una discontinuïtat evitable en $x=1$.»}
  \motiu{El límit que s'ha calculat és el de $x\to 1$ (val $\tfrac{1}{8}$), i
    $x=1$ és el zero del denominador que se simplifica.
    En $x=3$ la funció és contínua.}
\end{errada}

\begin{errada}{Exercici 102}
  \ondiu{$\dlim_{x\to -1^-} f(x)=5$}
  \hadedir{$\dlim_{x\to -1^-} f(x)=-5$}
  \motiu{Per a $x<-1$ la funció és $3x-2$, i $3\cdot(-1)-2=-5$.
    La conclusió no canvia: els límits laterals ($-5$ i $-4$) són
    diferents i en $x=-1$ hi ha una discontinuïtat de salt finit.}
\end{errada}

%\senseerrades{103, 40 i 106}

%% ============================================================
%  Setmana 4 — Unitat 7 · Límits i continuïtat
%  Exercicis: 112, 113, 114, 43, 120
% ============================================================

\setmana{4}{Unitat 7 \textperiodcentered{} Límits i continuïtat}{11 d'octubre de 2026}%
        {112, 113, 114, 43 i 120}

\begin{errada}{Exercici 43}
  \ondiu{$g(2)-4<0$}
  \hadedir{$g(2)=-4<0$}
  \motiu{Hi falta el signe igual: $g(2)=3^{-2+2}-5=1-5=-4$.
    El teorema de Bolzano necessita que $g$ tingui signes diferents als
    extrems de l'interval, i els té: $g(0)=4>0$ i $g(2)=-4<0$.}
\end{errada}

%\senseerrades{112, 113, 114 i 120}

% ---------- Observacions (només surten amb \obstrue) ----------
\ifobs
\bigskip
\noindent
Dues observacions que no són errades, però que poden confondre:

\begin{observacio}{Exercici 113}
  El valor $x=-3$ és una arrel exacta, perquè
  $x^3+6x^2+12x+9=(x+2)^3+1$. Si proves amb $x=-3$ trobaràs $f(-3)=0$ en lloc
  d'un canvi de signe: no t'has equivocat, has trobat l'arrel. Qualsevol
  interval de longitud $2$ que la contingui, com ara $(-4,-2)$, respon la
  pregunta.
\end{observacio}

\begin{observacio}{Exercici 120}
  L'enunciat anomena $f(x)=x+1$, i la solució torna a fer servir la lletra $f$
  per a una funció diferent: la resta de les dues funcions que es tallen (per
  exemple, $f(x)=x+1-e^{x-1}$ a l'apartat a). El punt on aquesta resta val
  zero és l'abscissa del punt de tall. D'altra banda, l'enunciat només demana
  demostrar que hi ha un punt de tall; la solució va més enllà i l'aproxima
  amb un error menor que una dècima.
\end{observacio}
\fi

%% ============================================================
%  Setmana 5 — Unitat 8 · Derivades
%  Exercicis: 37, 86, 87, 88, 92
% ============================================================

\setmana{5}{Unitat 8 \textperiodcentered{} Derivades}{18 d'octubre de 2026}%
        {37, 86, 87, 88 i 92}

\noindent
No hi ha errors.

%\senseerrades{37, 86, 87, 88 i 92}

% ---------- Observacions (només surten amb \obstrue) ----------
\ifobs
\bigskip
\noindent
Una observació, que no és cap errada:

\begin{observacio}{Exercici 87}
  L'enunciat demana calcular la funció derivada «utilitzant la definició»,
  però el solucionari només hi dona el resultat final de cada apartat, sense
  el desenvolupament. Els resultats són correctes. Si vols veure un model del
  desenvolupament pas a pas, mira la solució de l'exercici 88.
\end{observacio}
\fi

%% ============================================================
%  Setmana 6 — Unitat 8 · Derivades
%  Exercicis: 30, 31, 32, 21, 40, 41
% ============================================================

\setmana{6}{Unitat 8 \textperiodcentered{} Derivades}{25 d'octubre de 2026}%
        {30, 31, 32, 21, 40 i 41}

\noindent
No hi ha errors.

%\senseerrades{30, 31, 32, 21, 40 i 41}

%% ============================================================
%  Setmana 7 — Unitat 8 · Derivades
%  Exercicis: 46, 54, 58, 93, 97
% ============================================================

\setmana{7}{Unitat 8 \textperiodcentered{} Derivades}{1 de novembre de 2026}%
        {46, 54, 58, 93 i 97}

\begin{errada}{Exercici 58}
  \ondiu{«Els punts de tall de la funció amb els eixos són:»}
  \hadedir{«Els punts de tall de la recta tangent amb els eixos són:»}
  \motiu{Els punts que es calculen a continuació, $(0,7)$ i
    $\left(\tfrac{7}{2},0\right)$, són els talls de la recta tangent
    $y=-2x+7$, que és la que forma el triangle amb els eixos. La funció $f$
    talla els eixos en uns altres punts: $(0,-1)$ i $(-1,0)$.}
\end{errada}

%\senseerrades{46, 54, 93 i 97}

% ---------- Observacions (només surten amb \obstrue) ----------
\ifobs
\bigskip
\noindent
Dues observacions menors:\par\nopagebreak

\begin{observacio}{Exercici 54}
  L'enunciat escriu la recta com a $f(x)=2x-123$, amb la mateixa lletra $f$
  de la funció $f(x)=x^3-x+4$. Són dues coses diferents: la recta és
  $y=2x-123$, que és com l'escriu la solució.
\end{observacio}

\begin{observacio}{Exercici 58}
  A la solució hi ha una errada tipogràfica sense importància: hi diu
  «traingle» en lloc de «triangle».
\end{observacio}
\fi

%% ============================================================
%  Setmana 8 — Unitat 9 · Aplicacions de les derivades
%  Exercicis: 39, 41, 42, 44, 50, 51, 58
% ============================================================

\setmana{8}{Unitat 9 \textperiodcentered{} Aplicacions de les derivades}{8 de novembre de 2026}%
        {39, 41, 42, 44, 50, 51 i 58}

\begin{errada}{Exercici 41}
  \ondiu{«A $x=$, $f'(x)$ passa de positiva a negativa; per tant, és un màxim.»}
  \hadedir{«A $x=2$, $f'(x)$ passa de positiva a negativa; per tant, és un màxim.»}
  \motiu{Hi falta el valor de $x$. És l'últim dels quatre punts on la derivada
    val $0$ ($x=-4$, $x=-2$, $x=0$ i $x=2$), i a la gràfica $f'$ hi passa de
    positiva a negativa.}
\end{errada}

\begin{errada}{Exercici 42 d), valor de $y'(2)$}
  \ondiu{$y'(2)=80-60-20>0$}
  \hadedir{$y'(2)=80-60=20>0$}
  \motiu{$y'(2)=5\cdot 2^4-15\cdot 2^2=80-60=20$. Tal com està imprès,
    $80-60-20$ valdria $0$.}
\end{errada}

\begin{errada}{Exercici 42 d), extrems relatius}
  \ondiu{«Té màxims relatius a $x=-\sqrt{3}$, $x=\sqrt{3}$.»}
  \hadedir{«Té un màxim relatiu a $x=-\sqrt{3}$ i un mínim relatiu a $x=\sqrt{3}$.»}
  \motiu{A $x=-\sqrt{3}$ la funció passa de créixer a decréixer (màxim), però
    a $x=\sqrt{3}$ passa de decréixer a créixer: hi ha un mínim. A $x=0$ la
    derivada també s'anul·la, però no canvia de signe i no hi ha extrem.}
\end{errada}

\begin{errada}{Exercici 51}
  \ondiu{«El signe de $y$ és constant.»}
  \hadedir{«El signe de $y'$ és constant.»}
  \motiu{La que no s'anul·la mai i és contínua és la derivada
    $y'(x)=3x^2-4x+6$, i per això té signe constant (positiu, perquè
    $y'(0)=6>0$). La funció $y$ sí que canvia de signe: per exemple,
    $y(0)=-13$ i $y(3)=14$.}
\end{errada}

\begin{errada}{Exercici 58 b)}
  \ondiu{$y'(x)=0 \;\rightarrow\; a=\pm 2$}
  \hadedir{$y'(x)=0 \;\rightarrow\; x=\pm 2$}
  \motiu{El que es resol és $3x^2-12=0$, i la incògnita és $x$. El paràmetre
    $a$ ja s'ha trobat a l'apartat a) i val $-12$.}
\end{errada}

%\senseerrades{39, 44 i 50}

% ---------- Observacions (només surten amb \obstrue) ----------
\ifobs
\bigskip
\noindent
Tres observacions que no són errades de càlcul, però que poden confondre:\par\nopagebreak

\begin{observacio}{Exercici 39 b)}
  On diu «La derivada té un grau» vol dir «La derivada té grau u»: la derivada
  d'una paràbola és un polinomi de grau $1$, que s'anul·la en un sol punt, i
  per això la paràbola té, com a molt, un extrem.
\end{observacio}

\begin{observacio}{Exercici 41}
  L'enunciat diu «A partir de la derivada de $f'(x)$». Vol dir: a partir de la
  gràfica de la derivada $f'(x)$, que és la que es dona.
\end{observacio}

\begin{observacio}{Exercici 51}
  La primera frase de la solució enuncia la implicació al revés. El que es fa
  servir és: si $y'(x)>0$ per a tot $x$, llavors la funció és sempre creixent.
  La implicació contrària no és certa en general: $y=x^3$ és sempre creixent
  i, en canvi, $y'(0)=0$.
\end{observacio}
\fi

%% ============================================================
%  Setmana 9 — Unitat 9 · Aplicacions de les derivades
%  Exercicis: 30, 66, 67, 70
% ============================================================

\setmana{9}{Unitat 9 \textperiodcentered{} Aplicacions de les derivades}{15 de novembre de 2026}%
        {30, 66, 67 i 70}

\begin{errada}{Exercici 67 d)}
  \ondiu{$y''=-\dfrac{6}{x^2}$}
  \hadedir{$y''=-\dfrac{6}{x^4}$}
  \motiu{Com que $y'=\dfrac{x^3+2}{x^3}=1+2x^{-3}$, la derivada segona és
    $y''=-6x^{-4}$. La conclusió no canvia: $y''<0$ per a tot $x\neq 0$ i la
    funció és còncava a $(-\infty,0)$ i a $(0,+\infty)$.}
\end{errada}

\begin{errada}{Exercici 67 f)}
  \ondiu{«$y''>0$ a $(-\infty,-2)$ i a $(2,+\infty)$ $\rightarrow$ $y$ és
    convexa a $(-\infty,-2)$ i a $(2,+\infty)$.»}
  \hadedir{Aquestes dues línies sobren.}
  \motiu{El domini de la funció és $[-2,2]$, com diu la mateixa solució: fora
    d'aquest interval la funció no existeix i no té sentit parlar-ne de la
    curvatura. La resposta és només que $y$ és còncava a $(-2,2)$.}
\end{errada}

\begin{errada}{Exercici 67 h)}
  \ondiu{$y''(x)=\dfrac{2(x^2-x)-(2x-1)^2}{(x^2-x)}$}
  \hadedir{$y''(x)=\dfrac{2(x^2-x)-(2x-1)^2}{(x^2-x)^2}$}
  \motiu{En derivar el quocient, el denominador queda elevat al quadrat. El pas
    següent de la solució ja el porta bé, perquè $x^2(x-1)^2=(x^2-x)^2$.}
\end{errada}

\begin{errada}{Exercici 67 n)}
  \ondiu{$y''(x)>0$ a $\left(-\tfrac{\pi}{4},\tfrac{\pi}{4}\right)$ i a
    $\left(\tfrac{\pi}{4},\tfrac{\pi}{2}\right)$;\quad
    $y''(x)<0$ a $\left(-\tfrac{\pi}{2},-\tfrac{\pi}{4}\right)$}
  \hadedir{$y''(x)>0$ a $\left(-\tfrac{\pi}{4},\tfrac{\pi}{4}\right)$;\quad
    $y''(x)<0$ a $\left(-\tfrac{\pi}{2},-\tfrac{\pi}{4}\right)$ i a
    $\left(\tfrac{\pi}{4},\tfrac{\pi}{2}\right)$}
  \motiu{$y''=2\cos 2x$. A l'interval
    $\left(\tfrac{\pi}{4},\tfrac{\pi}{2}\right)$ l'angle $2x$ és entre
    $\tfrac{\pi}{2}$ i $\pi$, on el cosinus és negatiu. Per tant, la funció és
    convexa a $\left(-\tfrac{\pi}{4},\tfrac{\pi}{4}\right)$ i còncava als
    altres dos intervals.}
\end{errada}

%\senseerrades{30, 66 i 70}

% ---------- Observacions (només surten amb \obstrue) ----------
\ifobs
\bigskip
\noindent
Una observació, que no és cap errada:\par\nopagebreak

\begin{observacio}{Exercici 66 c) i d)}
  Les respostes «Dos punts d'inflexió» i «$n-2$ punts d'inflexió» són el
  nombre màxim. Un polinomi de quart grau en pot tenir dos o cap: per exemple,
  $y=x^4$ no en té cap, perquè $y''=12x^2$ no canvia de signe.
\end{observacio}
\fi

%% ============================================================
%  Setmana 10 — Unitat 9 · Aplicacions de les derivades
%  Exercicis: 34, 77, 79, 86, 87
% ============================================================

\setmana{10}{Unitat 9 \textperiodcentered{} Aplicacions de les derivades}{22 de novembre de 2026}%
        {34, 77, 79, 86 i 87}

\noindent
No hi ha errors.

%\senseerrades{34, 77, 79, 86 i 87}

%% ============================================================
%  Setmana 11 — Unitat 10 · Representació de funcions
%  Exercicis: 37, 38, 43, 45, 62, 63
% ============================================================

\setmana{11}{Unitat 10 \textperiodcentered{} Representació de funcions}{29 de novembre de 2026}%
        {37, 38, 43, 45, 62 i 63}

\begin{errada}{Exercici 37}
  \ondiu{«$f'(x)>0$ en $(-\infty,-2)\cup(2,6)$» \; i \;
    «$f'(x)<0$ en $(-2,2)\cup(6,+\infty)$»}
  \hadedir{«$f'(x)>0$ en $(-\infty,x_1)\cup(x_2,x_3)$» \; i \;
    «$f'(x)<0$ en $(x_1,x_2)\cup(x_3,+\infty)$»}
  \motiu{El signe de $f'$ canvia als punts on la gràfica talla l'eix $X$, que
    l'enunciat anomena $x_1$, $x_2$ i $x_3$ i que la mateixa solució fa servir
    a la línia següent. A la gràfica, $x_1$ és entre $-2$ i $-1$ i $x_3$ entre
    $5$ i $6$: no són $-2$ ni $6$.}
\end{errada}

\begin{errada}{Exercici 38 (enunciat)}
  \ondiu{«Creixement: $(-\infty,0)\cup(2\;-\infty)$.»}
  \hadedir{«Creixement: $(-\infty,0)\cup(2,+\infty)$.»}
  \motiu{És una errada de l'enunciat. La funció creix abans del màxim $(0,0)$
    i després del mínim $(2,4)$, com es veu a la gràfica de la solució.}
\end{errada}

\begin{errada}{Exercici 45 b)}
  \ondiu{$\mathrm{Dom}\,f=\R-\{3\}$}
  \hadedir{$\mathrm{Dom}\,f=\R-\{-3\}$}
  \motiu{El denominador és $x+3$, que s'anul·la en $x=-3$.}
\end{errada}

\begin{errada}{Exercici 62 c)}
  \ondiu{$\dlim_{x\to-\infty}\,(x^4+2x^3-2x-1)=-\infty$}
  \hadedir{$\dlim_{x\to-\infty}\,(x^4+2x^3-2x-1)=+\infty$}
  \motiu{El terme dominant és $x^4$, que és positiu tant si $x\to-\infty$ com
    si $x\to+\infty$.}
\end{errada}

\begin{errada}{Exercici 62 d)}
  \ondiu{$\dlim_{x\to+\infty}\,(-2x^4+3x^2-x-1)=+\infty$}
  \hadedir{$\dlim_{x\to+\infty}\,(-2x^4+3x^2-x-1)=-\infty$}
  \motiu{El terme dominant és $-2x^4$, que tendeix a $-\infty$ tant si
    $x\to-\infty$ com si $x\to+\infty$.}
\end{errada}

\begin{errada}{Exercici 63 c)}
  \ondiu{$\dlim_{x\to-\infty}\dfrac{2x}{x+1}=2$ \; i \;
    $\dlim_{x\to+\infty}\dfrac{2x}{x+1}=2$}
  \hadedir{$\dlim_{x\to-\infty}\dfrac{4x^2}{x^2-9}=4$ \; i \;
    $\dlim_{x\to+\infty}\dfrac{4x^2}{x^2-9}=4$}
  \motiu{Els límits que hi ha impresos són d'una altra funció. Amb la de
    l'apartat, $y=\dfrac{4x^2}{x^2-9}$, el límit és $4$, que és el que dona
    l'asímptota horitzontal $y=4$ de la mateixa solució.}
\end{errada}

%\senseerrades{43}

%% ============================================================
%  Setmana 12 — Unitat 10 · Representació de funcions
%  Exercicis: 75, 78, 88, 91, 41, 100
% ============================================================

\setmana{12}{Unitat 10 \textperiodcentered{} Representació de funcions}{6 de desembre de 2026}%
        {75, 78, 88, 91, 41 i 100}

\begin{errada}{Exercici 75 a), signe de la derivada}
  \ondiu{«En $(-5,1)$, $y'>0$ $\rightarrow$ Funció decreixent»}
  \hadedir{«En $(-5,1)$, $y'<0$ $\rightarrow$ Funció decreixent»}
  \motiu{Entre les dues arrels la derivada és negativa: per exemple,
    $y'(0)=-15$.}
\end{errada}

\begin{errada}{Exercici 75 a), màxim}
  \ondiu{$x=-5 \;\rightarrow\; y=203 \;\rightarrow\; (-5,203)$ màxim}
  \hadedir{$x=-5 \;\rightarrow\; y=103 \;\rightarrow\; (-5,103)$ màxim}
  \motiu{$y(-5)=(-5)^3+6\cdot(-5)^2-15\cdot(-5)+3=-125+150+75+3=103$.}
\end{errada}

\begin{errada}{Exercici 75 b)}
  \ondiu{$x=-\tfrac{1}{2} \rightarrow y=-6 \rightarrow
    \left(-\tfrac{1}{2},-6\right)$ màxim;\quad
    $x=-\tfrac{1}{2} \rightarrow y=6 \rightarrow
    \left(-\tfrac{1}{2},6\right)$ mínim}
  \hadedir{$x=-\tfrac{1}{2} \rightarrow y=-4 \rightarrow
    \left(-\tfrac{1}{2},-4\right)$ màxim;\quad
    $x=\tfrac{1}{2} \rightarrow y=4 \rightarrow
    \left(\tfrac{1}{2},4\right)$ mínim}
  \motiu{El numerador val $4\cdot\tfrac{1}{4}+1=2$ en tots dos punts, i
    $2:\tfrac{1}{2}=4$. Per tant, $y\left(\tfrac{1}{2}\right)=4$ i
    $y\left(-\tfrac{1}{2}\right)=-4$. A més, el mínim és a $x=\tfrac{1}{2}$,
    no a $x=-\tfrac{1}{2}$.}
\end{errada}

\begin{errada}{Exercici 75 d)}
  \ondiu{«$y'=4x^3-72x^2=0 \rightarrow x=0$. En $\R-\{0\}$, $y'<0$
    $\rightarrow$ Funció decreixent»}
  \hadedir{«$y'=4x^3-72x^2=0 \rightarrow x=0$, $x=18$.
    En $(-\infty,0)\cup(0,18)$, $y'<0$ $\rightarrow$ Funció decreixent.
    En $(18,+\infty)$, $y'>0$ $\rightarrow$ Funció creixent.
    $x=18 \rightarrow y=-34\,992 \rightarrow (18,-34\,992)$ mínim»}
  \motiu{$y'=4x^2(x-18)$ també s'anul·la en $x=18$, i allà passa de negativa a
    positiva. En $x=0$ la derivada val $0$ però no canvia de signe: no hi ha
    extrem.}
\end{errada}

\begin{errada}{Exercici 78 b), domini}
  \ondiu{$\mathrm{Dom}=\R$}
  \hadedir{$\mathrm{Dom}=\R-\{-2\}$}
  \motiu{El denominador $x+2$ s'anul·la en $x=-2$.}
\end{errada}

\begin{errada}{Exercici 78 b), interval}
  \ondiu{«En $(-\infty,2)$, $y''>0$ $\rightarrow$ Funció convexa»}
  \hadedir{«En $(-\infty,-2)$, $y''>0$ $\rightarrow$ Funció convexa»}
  \motiu{$y''=\dfrac{-8}{(x+2)^3}$ canvia de signe en $x=-2$: és positiva per a
    $x<-2$ i negativa per a $x>-2$.}
\end{errada}

\begin{errada}{Exercici 78 d), derivada segona}
  \ondiu{$y''=\dfrac{12x^2+4x}{(x^2-1)^3}\neq 0$}
  \hadedir{$y''=\dfrac{12x^2+4}{(x^2-1)^3}\neq 0$}
  \motiu{El numerador és $12x^2+4$, que no s'anul·la mai. Amb $4x$ s'anul·laria
    en $x=0$, i la solució no podria dir que $y''\neq 0$.}
\end{errada}

\begin{errada}{Exercici 78 d), interval}
  \ondiu{«En $(1,1)$, $y''<0$ $\rightarrow$ Funció còncava»}
  \hadedir{«En $(-1,1)$, $y''<0$ $\rightarrow$ Funció còncava»}
  \motiu{Entre $-1$ i $1$ el denominador $(x^2-1)^3$ és negatiu i el numerador
    és positiu.}
\end{errada}

\begin{errada}{Exercici 41}
  \ondiu{«$f'(x)=\dfrac{-x^2-1}{(x^2-1)^2}<0 \;\rightarrow\; f(x)$ creixent»}
  \hadedir{«$f'(x)=\dfrac{-x^2-1}{(x^2-1)^2}<0 \;\rightarrow\; f(x)$ decreixent»}
  \motiu{Si la derivada és negativa, la funció decreix. És el que es veu a la
    gràfica de la mateixa solució, on les tres branques baixen.}
\end{errada}

\begin{errada}{Exercici 100 a)}
  \ondiu{$y''=\dfrac{-1}{2(2-x)\sqrt{2-x}}<0$}
  \hadedir{$y''=\dfrac{-1}{4(2-x)\sqrt{2-x}}<0$}
  \motiu{De $y'=-\tfrac{1}{2}(2-x)^{-1/2}$ surt $y''=-\tfrac{1}{4}(2-x)^{-3/2}$.
    La conclusió (funció còncava) no canvia.}
\end{errada}

\begin{errada}{Exercici 100 c), límit}
  \ondiu{$\dlim_{x\to-\infty}\dfrac{\sqrt{x^2-9}}{x}=1 \;\rightarrow\; m=-1$}
  \hadedir{$\dlim_{x\to-\infty}\dfrac{\sqrt{x^2-9}}{x}=-1 \;\rightarrow\; m=-1$}
  \motiu{Quan $x\to-\infty$, el numerador és positiu i el denominador negatiu:
    el límit és $-1$, que és el pendent $m$ que la solució fa servir.}
\end{errada}

\begin{errada}{Exercici 100 c), títol}
  \ondiu{«Creixement i decreixement:» (a sobre de la derivada segona)}
  \hadedir{«Concavitat i convexitat:»}
  \motiu{El títol «Creixement i decreixement» hi surt dues vegades. La segona,
    la que porta $y''$, és l'estudi de la curvatura.}
\end{errada}

\begin{errada}{Exercici 100 d)}
  \ondiu{$y''=\dfrac{-1}{4(x+3)\sqrt{x+3}}>0$}
  \hadedir{$y''=\dfrac{1}{4(x+3)\sqrt{x+3}}>0$}
  \motiu{De $y'=-\tfrac{1}{2}(x+3)^{-1/2}$ surt $y''=\tfrac{1}{4}(x+3)^{-3/2}$,
    que és positiva. La conclusió (funció convexa) és correcta.}
\end{errada}

%\senseerrades{88 i 91}

% ---------- Observacions (només surten amb \obstrue) ----------
\ifobs
\bigskip
\noindent
Una observació, que no és cap errada:\par\nopagebreak

\begin{observacio}{Exercici 75 e)}
  La solució estudia el creixement, però no diu l'extrem relatiu que demana
  l'enunciat: hi ha un mínim a $(0,1)$.
\end{observacio}
\fi

%% ============================================================
%  Setmana 13 — Unitat 13 · Probabilitat
%  Exercicis: 28, 29, 40, 60, 62, 32, 33, 35
% ============================================================

\setmana{13}{Unitat 13 \textperiodcentered{} Probabilitat}{13 de desembre de 2026}%
        {28, 29, 40, 60, 62, 32, 33 i 35}

\begin{errada}{Exercici 40 a)}
  \ondiu{$VR_{5,6}=5^6=15\,625$}
  \hadedir{$4\cdot VR_{5,5}=4\cdot 5^5=12\,500$}
  \motiu{Els dígits parells són $0$, $2$, $4$, $6$ i $8$, però un nombre de
    6 xifres no pot començar per $0$: per a la primera xifra només hi ha
    4 opcions. El resultat $15\,625$ compta també seqüències com $002468$, que
    no són nombres de 6 xifres.}
\end{errada}

\begin{errada}{Exercici 60 a)}
  \ondiu{$P(A\cap B)=1-P(\overline{A}\cap\overline{B})=1-0{,}3=0{,}7$}
  \hadedir{$P(A\cup B)=1-P(\overline{A}\cap\overline{B})=1-0{,}3=0{,}7$}
  \motiu{«No s'ha complert alguna de les dues normes» és la unió $A\cup B$,
    que és el contrari de $\overline{A}\cap\overline{B}$ (complir-les totes
    dues). La intersecció val una altra cosa:
    $P(A\cap B)=0{,}23+0{,}65-0{,}7=0{,}18$.}
\end{errada}

%\senseerrades{28, 29, 62, 32, 33 i 35}

% ---------- Observacions (només surten amb \obstrue) ----------
\ifobs
\bigskip
\noindent
Una observació menor:\par\nopagebreak

\begin{observacio}{Exercici 35}
  A la solució hi ha una errada tipogràfica sense importància: hi diu
  «probabiltat» en lloc de «probabilitat».
\end{observacio}
\fi

%% ============================================================
%  Setmana 14 — Unitat 13 · Probabilitat
%  Exercicis: 36, 96, 103, 104, 105, 106, 107, 108
% ============================================================

\setmana{14}{Unitat 13 \textperiodcentered{} Probabilitat}{20 de desembre de 2026}%
        {36, 96, 103, 104, 105, 106, 107 i 108}

\noindent
No hi ha errors.

%\senseerrades{36, 96, 103, 104, 105, 106, 107 i 108}

% ---------- Observacions (només surten amb \obstrue) ----------
\ifobs
\bigskip
\noindent
Una observació menor:\par\nopagebreak

\begin{observacio}{Exercicis 36 i 96}
  A les solucions hi ha dues errades tipogràfiques sense importància:
  «probabiltat» en lloc de «probabilitat» (exercici 36) i «perfeccionamient»
  en lloc de «perfeccionament» (exercici 96).
\end{observacio}
\fi

%% ============================================================
%  Setmana 15 — Unitat 14 · Distribucions de probabilitat
%  Exercicis: 43, 46, 44, 50
% ============================================================

\setmana{15}{Unitat 14 \textperiodcentered{} Distribucions de probabilitat}{27 de desembre de 2026}%
        {43, 46, 44 i 50}

\noindent
No hi ha errors.

%\senseerrades{43, 46, 44 i 50}

% ---------- Observacions (només surten amb \obstrue) ----------
\ifobs
\bigskip
\noindent
Una observació menor:\par\nopagebreak

\begin{observacio}{Exercicis 46 i 50}
  A les solucions hi ha dues errades tipogràfiques sense importància:
  «propabilitat» en lloc de «probabilitat» (exercici 46 a) i, a l'exercici 50,
  un parèntesi que no es tanca: ha de dir
  $1-\bigl(P(X=0)+P(X=1)+P(X=2)\bigr)$.
\end{observacio}
\fi

%% ============================================================
%  Setmana 16 — Unitat 14 · Distribucions de probabilitat
%  Exercicis: 21, 53, 100, 101
% ============================================================

\setmana{16}{Unitat 14 \textperiodcentered{} Distribucions de probabilitat}{3 de gener de 2027}%
        {21, 53, 100 i 101}

\begin{errada}{Exercici 53 b)}
  \ondiu{$P(X=3)=0{,}188$}
  \hadedir{$P(X=12)=0{,}000000007$}
  \motiu{$X$ compta els enquestats que responen «NS/NC». Si exactament 3 no
    responen «NS/NC», els altres 12 sí que ho fan: cal
    $P(X=12)=\binom{15}{12}\cdot 0{,}13^{12}\cdot 0{,}87^{3}\approx 7\cdot 10^{-9}$.
    El valor $0{,}188$ és la probabilitat que exactament 3 responguin «NS/NC»:
    seria la resposta si a l'enunciat hi sobrés el «no».}
\end{errada}

%\senseerrades{21, 100 i 101}

%% ============================================================
%  Setmana 17 — Unitat 10 · Representació de funcions
%  Exercicis: 84, 108, 123, 124
% ============================================================

\setmana{17}{Unitat 10 \textperiodcentered{} Representació de funcions}{10 de gener de 2027}%
        {84, 108, 123 i 124}

\begin{errada}{Exercici 84 d)}
  \ondiu{$f\left(\tfrac{2}{3}\right)=\dfrac{26}{27} \;\rightarrow\;
    \left(\tfrac{2}{3},\tfrac{26}{27}\right)$ punt d'inflexió}
  \hadedir{$f\left(\tfrac{2}{3}\right)=\dfrac{20}{27} \;\rightarrow\;
    \left(\tfrac{2}{3},\tfrac{20}{27}\right)$ punt d'inflexió}
  \motiu{$f\left(\tfrac{2}{3}\right)=\tfrac{8}{27}-\tfrac{24}{27}-\tfrac{18}{27}
    +\tfrac{54}{27}=\tfrac{20}{27}$.}
\end{errada}

\begin{errada}{Exercici 84 g)}
  \ondiu{«$f''(x)=-2>0$ \; $\forall x\in\R$»}
  \hadedir{«$f''(x)=-2<0$ \; $\forall x\in\R$»}
  \motiu{$-2$ és negatiu, i per això la funció és còncava, com diu la mateixa
    solució.}
\end{errada}

\begin{errada}{Exercici 108 b)}
  \ondiu{«Eix $Y$: $f(0)=0 \;\rightarrow\; (0,6)$»}
  \hadedir{«Eix $Y$: $f(0)=6 \;\rightarrow\; (0,6)$»}
  \motiu{Per a $x>-2$ la funció és $6-x^2$, i $f(0)=6$. El punt $(0,6)$ ja és
    correcte.}
\end{errada}

\begin{errada}{Exercici 108 c), continuïtat}
  \ondiu{$\dlim_{x\to-2^-}f(x)=\dlim_{x\to 2^-}(x^3-1)=7$}
  \hadedir{$\dlim_{x\to 2^-}f(x)=\dlim_{x\to 2^-}(x^3-1)=7$}
  \motiu{El punt que s'estudia és $x=2$, on canvia l'expressió de la funció.}
\end{errada}

\begin{errada}{Exercici 108 c), derivada segona}
  \ondiu{$\dfrac{14}{(x-1)^3}$ \; si \; $x<2$}
  \hadedir{$\dfrac{14}{(x-1)^3}$ \; si \; $x>2$}
  \motiu{És la derivada segona del tercer tros, $\dfrac{7}{x-1}$, que és el de
    $x>2$.}
\end{errada}

\begin{errada}{Exercici 108 d)}
  \ondiu{«$2-3\cos x=0 \;\rightarrow\; x=\arccos\left(\tfrac{2}{3}\right)
    \;\rightarrow\; \left(\arccos\left(\tfrac{2}{3}\right),0\right)$»}
  \hadedir{«$2-3\cos x=0 \;\rightarrow\;
    x=\pm\arccos\left(\tfrac{2}{3}\right)+2k\pi$, amb $x\leq 0$. El tall més
    proper a l'origen és $\left(-\arccos\left(\tfrac{2}{3}\right),0\right)$»}
  \motiu{El tros $2-3\cos x$ només val per a $x\leq 0$, i
    $\arccos\left(\tfrac{2}{3}\right)\approx 0{,}84$ és positiu: allà la funció
    és $2^{-x+1}-3$ i no val $0$. Els talls amb l'eix $X$ són tots a
    l'esquerra de l'origen i n'hi ha infinits, com es veu a la gràfica de la
    solució.}
\end{errada}

%\senseerrades{123 i 124}

% ---------- Observacions (només surten amb \obstrue) ----------
\ifobs
\bigskip
\noindent
Tres observacions que no són errades de càlcul:\par\nopagebreak

\begin{observacio}{Exercici 84 d)}
  Al segon extrem,
  $\left(\tfrac{2+\sqrt{7}}{3},\tfrac{20-14\sqrt{7}}{27}\right)$, la solució
  no diu de quin tipus és: és un mínim.
\end{observacio}

\begin{observacio}{Exercici 108 b) i c)}
  La solució només hi dona els mínims. També hi ha màxims relatius: a b), el
  punt $(0,6)$; a c), el punt $(2,7)$, on la funció passa de créixer a
  decréixer.
\end{observacio}

\begin{observacio}{Exercici 124 c)}
  La solució escriu $T'(x)$ i $T(x)$ en lloc de $T'(t)$ i $T(t)$, i dona
  l'interval $(-\infty,2)$ tot i que el temps comença a $t=0$.
\end{observacio}
\fi

%% ============================================================
%  Setmana 18 — Unitat 1 · Matrius
%  Exercicis: 12, 13, 15, 16, 17
% ============================================================

\setmana{18}{Unitat 1 \textperiodcentered{} Matrius}{17 de gener de 2027}%
        {12, 13, 15, 16 i 17}

\noindent
No hi ha errors.

%\senseerrades{12, 13, 15, 16 i 17}

%% ============================================================
%  Setmana 19 — Unitat 1 · Matrius
%  Exercicis: 42, 49, 69, 81, 85
% ============================================================

\setmana{19}{Unitat 1 \textperiodcentered{} Matrius}{24 de gener de 2027}%
        {42, 49, 69, 81 i 85}

\begin{errada}{Exercici 49 d)}
  \ondiu{$BC=\begin{pmatrix}-7 & 3 & 3\\ -1 & -1 & 8\end{pmatrix}$}
  \hadedir{$BC=\begin{pmatrix}-8 & 3 & 3\\ -4 & -1 & 8\end{pmatrix}$}
  \motiu{Els dos elements de la primera columna estan malament:
    $(-2)\cdot 4+1\cdot 0+1\cdot 0=-8$ i $(-1)\cdot 4+0\cdot 0+3\cdot 0=-4$.}
\end{errada}

%\senseerrades{42, 69, 81 i 85}

% ---------- Observacions (només surten amb \obstrue) ----------
\ifobs
\bigskip
\noindent
Dues observacions, que no són cap errada:\par\nopagebreak

\begin{observacio}{Exercici 42}
  L'enunciat demana dos exemples de cada tipus de matriu i la solució només
  en dona un.
\end{observacio}

\begin{observacio}{Exercici 85 b)}
  La solució no exclou $a=0$ ni $b=0$. Si $a=b=0$, $B$ és la matriu nul·la, que no és nilpotent de grau 2 perquè ja és nul·la amb l'exponent 1. Cal afegir-hi $a\neq 0$ a la primera família de matrius i $b\neq 0$ a la segona.
\end{observacio}
\fi

%% ============================================================
%  Setmana 20 — Unitat 2 · Determinants
%  Exercicis: 35, 36, 37, 41, 44
% ============================================================

\setmana{20}{Unitat 2 \textperiodcentered{} Determinants}{31 de gener de 2027}%
        {35, 36, 37, 41 i 44}

\noindent
No hi ha errors.

%\senseerrades{35, 36, 37, 41 i 44}

% ---------- Observacions (només surten amb \obstrue) ----------
\ifobs
\bigskip
\noindent
Una observació menor:\par\nopagebreak

\begin{observacio}{Exercici 44}
  A l'enunciat hi ha una paraula repetida: «perquè que el determinant» ha de
  ser «perquè el determinant».
\end{observacio}
\fi

%% ============================================================
%  Setmana 21 — Unitat 2 · Determinants
%  Exercicis: 79, 80, 83, 87
% ============================================================

\setmana{21}{Unitat 2 \textperiodcentered{} Determinants}{7 de febrer de 2027}%
        {79, 80, 83 i 87}

\begin{errada}{Exercici 80 b)}
  \ondiu{$\begin{vmatrix}3&-5&1\\4&1&-2\\7&-4&-2\end{vmatrix}=0$}
  \hadedir{$\begin{vmatrix}3&-5&1\\4&1&-2\\7&-4&-1\end{vmatrix}=0$}
  \motiu{La fila que s'ha afegit a la matriu és $(7,-4,-1)$, la suma de les
    altres dues, i per això el determinant és $0$. Amb $-2$ al final, el
    determinant valdria $-23$.}
\end{errada}

%\senseerrades{79, 83 i 87}

% ---------- Observacions (només surten amb \obstrue) ----------
\ifobs
\bigskip
\noindent
Una observació, que no és cap errada:\par\nopagebreak

\begin{observacio}{Exercici 79 d)}
  El rang 2 és correcte, però la justificació és incompleta. Dels dos determinants d'ordre 3 que es calculen, només el segon és un menor orlat del menor $\begin{vmatrix}1&3\\0&5\end{vmatrix}$ que s'agafa al final. N'hi falta l'altre, $\begin{vmatrix}2&1&3\\-2&0&5\\2&3&19\end{vmatrix}$, que també val $0$.
\end{observacio}
\fi

%% ============================================================
%  Setmana 22 — Unitat 2 · Determinants
%  Exercicis: 95, 96, 99, 100
% ============================================================

\setmana{22}{Unitat 2 \textperiodcentered{} Determinants}{14 de febrer de 2027}%
        {95, 96, 99 i 100}

\begin{errada}{Exercici 95 (enunciat)}
  \ondiu{III) $-\dfrac{1}{2}\begin{pmatrix}4&-3\\-2&2\end{pmatrix}$}
  \hadedir{III) $-\dfrac{1}{2}\begin{pmatrix}4&-3\\-2&1\end{pmatrix}$}
  \motiu{És una errada de l'enunciat. La inversa de
    $\begin{pmatrix}1&3\\2&4\end{pmatrix}$ és
    $\dfrac{1}{-2}\begin{pmatrix}4&-3\\-2&1\end{pmatrix}$: l'últim element és
    $1$. Tal com està impresa, l'opció III) no és la inversa de cap de les
    tres matrius; la relació b) $\rightarrow$ III) de la solució és la que es
    volia.}
\end{errada}

%\senseerrades{96, 99 i 100}

%% ============================================================
%  Setmana 23 — Unitat 3 · Sistemes d'equacions
%  Exercicis: 38, 41, 42, 39, 43
% ============================================================

\setmana{23}{Unitat 3 \textperiodcentered{} Sistemes d'equacions}{21 de febrer de 2027}%
        {38, 41, 42, 39 i 43}

\begin{errada}{Exercici 41 d)}
  \ondiu{$2x-z=2$}
  \hadedir{$2x+z=2$}
  \motiu{La segona fila de la matriu és $(2\;\;0\;\;1)$: el coeficient de $z$
    és $1$.}
\end{errada}

\begin{errada}{Exercici 43 f)}
  \ondiu{$F_3=F_3-2F_1$}
  \hadedir{$F_3=F_3-2F_2$}
  \motiu{L'operació que anul·la la tercera fila, $(0\;\;10\;\;34\mid 2)$, és
    restar-li el doble de la segona, $(0\;\;5\;\;17\mid 1)$. La matriu que
    surt a continuació ja és la correcta.}
\end{errada}

%\senseerrades{38, 42 i 39}

%% ============================================================
%  Setmana 24 — Unitat 3 · Sistemes d'equacions
%  Exercicis: 55, 56
% ============================================================

\setmana{24}{Unitat 3 \textperiodcentered{} Sistemes d'equacions}{28 de febrer de 2027}%
        {55 i 56}

\begin{errada}{Exercici 55 d)}
  \ondiu{$\begin{vmatrix}1&1\\0&-1\end{vmatrix}=1\neq 0$}
  \hadedir{$\begin{vmatrix}1&1\\0&1\end{vmatrix}=1\neq 0$}
  \motiu{El menor es pren de la matriu $A$, i l'element de la segona fila és
    $1$. Amb $-1$ el determinant valdria $-1$.}
\end{errada}

\begin{errada}{Exercici 56 d)}
  \ondiu{$\begin{vmatrix}3&2&7\\1&-1&6\\6&-1&-3\end{vmatrix}=110\neq 0$}
  \hadedir{$\begin{vmatrix}3&2&7\\1&-1&6\\6&-1&3\end{vmatrix}=110\neq 0$}
  \motiu{L'última columna és la dels termes independents, $7$, $6$ i $3$.
    Amb $3$ el determinant val $110$; amb $-3$ valdria $140$.}
\end{errada}

%% ============================================================
%  Setmana 25 — Unitat 3 · Sistemes d'equacions
%  Exercicis: 60, 64, 66
% ============================================================

\setmana{25}{Unitat 3 \textperiodcentered{} Sistemes d'equacions}{7 de març de 2027}%
        {60, 64 i 66}

\begin{errada}{Exercici 60 a)}
  \ondiu{«Si $m\neq 3$, $m\neq -3$ $\rightarrow$
    $\mathrm{Rang}(A)=\mathrm{Rang}(A^*)=2<$ nre. d'incògnites $\rightarrow$
    Sistema compatible determinat»}
  \hadedir{«Si $m\neq 3$, $m\neq -3$ $\rightarrow$
    $\mathrm{Rang}(A)=\mathrm{Rang}(A^*)=2=$ nre. d'incògnites $\rightarrow$
    Sistema compatible determinat»}
  \motiu{El sistema té dues incògnites, $x$ i $y$, i el rang és $2$: són
    iguals, i per això és compatible determinat. Si el rang fos més petit que
    el nombre d'incògnites, seria indeterminat.}
\end{errada}

\begin{errada}{Exercici 60 d)}
  \ondiu{«Si $m=-4$ $\rightarrow$ $\mathrm{Rang}(A)=1$ $\rightarrow$
    $\mathrm{Rang}(A)\neq\mathrm{Rang}(A)=2$»}
  \hadedir{«Si $m=-4$ $\rightarrow$ $\mathrm{Rang}(A)=1$ $\rightarrow$
    $\mathrm{Rang}(A)\neq\mathrm{Rang}(A^*)=2$»}
  \motiu{La que té rang $2$ és la matriu ampliada $A^*$; hi falta l'asterisc.}
\end{errada}

\begin{errada}{Exercici 60 f), cas $m=0$}
  \ondiu{$\begin{vmatrix}-1&2\\1&2\end{vmatrix}=-6\neq 0$}
  \hadedir{$\begin{vmatrix}-1&2\\2&2\end{vmatrix}=-6\neq 0$}
  \motiu{El menor surt de les files 2 i 3 de $A$ (columnes 1 i 3), que per a
    $m=0$ són $(-1\;\;0\;\;2)$ i $(2\;\;0\;\;2)$. Amb un $1$ el determinant
    valdria $-4$. Al cas $m=3$ la solució l'escriu bé.}
\end{errada}

%\senseerrades{64 i 66}

% ---------- Observacions (només surten amb \obstrue) ----------
\ifobs
\bigskip
\noindent
Dues observacions que no són errades de càlcul:\par\nopagebreak

\begin{observacio}{Exercici 64}
  L'enunciat demana resoldre el sistema quan és compatible indeterminat
  ($a=1$) i la solució només el discuteix. Per a $a=1$ el sistema és
  $2x+y+z=2$, $x+y=1$, $-y+z=0$, i les solucions són $x=1-\lambda$,
  $y=\lambda$, $z=\lambda$ amb $\lambda\in\R$.
\end{observacio}

\begin{observacio}{Exercici 60 d) i e)}
  Dues errades tipogràfiques sense importància: a d) hi diu «Rango ($A$)» en
  lloc de «Rang ($A$)», i a e), a «Si $m\neq 1$, $m$ $\rightarrow$», hi sobra
  la segona $m$.
\end{observacio}
\fi

%% ============================================================
%  Setmana 26 — Unitat 3 · Sistemes d'equacions
%  Exercicis: 90, 92, 94, 100
% ============================================================

\setmana{26}{Unitat 3 \textperiodcentered{} Sistemes d'equacions}{14 de març de 2027}%
        {90, 92, 94 i 100}

\begin{errada}{Exercici 90}
  \ondiu{$F_2=2F_1-F_2$ \; (primera operació)}
  \hadedir{$F_1=2F_1-F_2$}
  \motiu{La fila que canvia és la primera: $2F_1-F_2=(10\;\;10\;\;0\mid 42)$,
    i les altres dues queden igual.}
\end{errada}

\begin{errada}{Exercici 100}
  \ondiu{$A=\begin{pmatrix}0&2&3\\4&0&3\\5&2&1\end{pmatrix}$ \; (i el mateix
    a la matriu $A^*$)}
  \hadedir{$A=\begin{pmatrix}0&2&3\\4&0&1\\5&2&0\end{pmatrix}$}
  \motiu{Les files són els coeficients de $2p_2+3p_3=440$, $4p_1+p_3=240$ i
    $5p_1+2p_2=360$. Els determinants de Cramer que venen després ja fan
    servir la matriu correcta.}
\end{errada}

%\senseerrades{92 i 94}

%% ============================================================
%  Setmana 27 — Unitat 5 · Rectes i plans en l'espai
%  Exercicis: 43, 45, 46, 47, 49
% ============================================================

\setmana{27}{Unitat 5 \textperiodcentered{} Rectes i plans en l'espai}{21 de març de 2027}%
        {43, 45, 46, 47 i 49}

\begin{errada}{Exercici 46}
  \ondiu{$\dfrac{x-1}{1}=\dfrac{y-2}{2}$ \; (equació implícita)}
  \hadedir{$\dfrac{x-1}{-1}=\dfrac{y-2}{2}$}
  \motiu{El denominador és la primera coordenada del vector director, $-1$.
    Amb $-1$ surt $2x+y-4=0$, que és el resultat de la solució; amb $1$
    sortiria $2x-y=0$.}
\end{errada}

\begin{errada}{Exercici 49 a)}
  \ondiu{$\dfrac{x=0}{2}=\dfrac{y-3}{5}=\dfrac{z-0}{1}$}
  \hadedir{$\dfrac{x-0}{2}=\dfrac{y-3}{5}=\dfrac{z-0}{1}$}
  \motiu{Al primer numerador hi ha d'haver un signe menys, no un igual.}
\end{errada}

\begin{errada}{Exercici 49 b)}
  \ondiu{$x-2\lambda$; \; $y=\lambda$; \; $z=\lambda$}
  \hadedir{$x=2\lambda$; \; $y=\lambda$; \; $z=\lambda$}
  \motiu{De $x-2y=0$ amb $y=\lambda$ surt $x=2\lambda$.}
\end{errada}

%\senseerrades{43, 45 i 47}

%% ============================================================
%  Setmana 28 — Unitat 5 · Rectes i plans en l'espai
%  Exercicis: 53, 54, 55, 56, 59
% ============================================================

\setmana{28}{Unitat 5 \textperiodcentered{} Rectes i plans en l'espai}{28 de març de 2027}%
        {53, 54, 55, 56 i 59}

\begin{errada}{Exercici 56}
  \ondiu{$z=\lambda$; \; $y=\lambda$; \; $z=\lambda$}
  \hadedir{$x=\lambda$; \; $y=\lambda$; \; $z=\lambda$}
  \motiu{La primera equació paramètrica és la de $x$.}
\end{errada}

\begin{errada}{Exercici 59 b)}
  \ondiu{$-z+6=0 \;\rightarrow\; z=-6 \;\rightarrow\; P_Z(0,0,6)$}
  \hadedir{$-z+6=0 \;\rightarrow\; z=6 \;\rightarrow\; P_Z(0,0,6)$}
  \motiu{De $-z+6=0$ surt $z=6$. El punt $P_Z(0,0,6)$ ja és correcte.}
\end{errada}

%\senseerrades{53, 54 i 55}

%% ============================================================
%  Setmana 29 — Unitat 5 · Rectes i plans en l'espai
%  Exercicis: 63, 66, 71, 72, 73
% ============================================================

\setmana{29}{Unitat 5 \textperiodcentered{} Rectes i plans en l'espai}{4 d'abril de 2027}%
        {63, 66, 71, 72 i 73}

\begin{errada}{Exercici 63}
  \ondiu{$\dfrac{x-4}{2}=\dfrac{y+2}{4}=\dfrac{z-1}{-2}$}
  \hadedir{$\dfrac{x-4}{-2}=\dfrac{y+2}{4}=\dfrac{z-1}{-2}$}
  \motiu{El vector director és $\overrightarrow{AB}=(-2,4,-2)$. A la línia
    següent, quan s'hi substitueix el punt $C$, la solució ja hi posa $-2$.}
\end{errada}

%\senseerrades{66, 71, 72 i 73}

%% ============================================================
%  Setmana 30 — Unitat 5 · Rectes i plans en l'espai
%  Exercicis: 74, 76, 77, 79
% ============================================================

\setmana{30}{Unitat 5 \textperiodcentered{} Rectes i plans en l'espai}{11 d'abril de 2027}%
        {74, 76, 77 i 79}

\begin{errada}{Exercici 77 b)}
  \ondiu{$M^*=\begin{pmatrix}1&-2&-1&-1\\2&1&5&5\\7&-4&7&7\end{pmatrix}$}
  \hadedir{$M^*=\begin{pmatrix}1&-2&3&-1\\2&1&-1&5\\7&-4&7&7\end{pmatrix}$}
  \motiu{La tercera columna és la dels coeficients de $z$ ($3$, $-1$ i $7$);
    a la solució hi surt repetida la dels termes independents.}
\end{errada}

%\senseerrades{74, 76 i 79}

% ---------- Observacions (només surten amb \obstrue) ----------
\ifobs
\bigskip
\noindent
Una observació, que no és cap errada de càlcul:\par\nopagebreak

\begin{observacio}{Exercicis 76 b) i 77}
  La solució escriu «$M^*=(\ldots)=2$» i «$M^*=(\ldots)=3$». Vol dir que el
  rang de $M^*$ és $2$ o $3$: una matriu no és igual a un nombre.
\end{observacio}
\fi

%% ============================================================
%  Setmana 31 — Unitat 5 · Rectes i plans en l'espai
%  Exercicis: 80, 81, 84, 85
% ============================================================

\setmana{31}{Unitat 5 \textperiodcentered{} Rectes i plans en l'espai}{18 d'abril de 2027}%
        {80, 81, 84 i 85}

\begin{errada}{Exercici 80}
  \ondiu{«Les rectes són paral·leles quan $m=2$.»}
  \hadedir{«Les rectes són paral·leles quan $m=-2$.»}
  \motiu{És el valor que surt del sistema de la línia anterior. Amb $m=2$ es
    compleix $m^2=4$, però no $2-m\neq 0$: les dues rectes serien la mateixa.
    El pla que es calcula després ja fa servir $m=-2$.}
\end{errada}

%\senseerrades{81, 84 i 85}

%% ============================================================
%  Setmana 32 — Unitat 6 · Angles i distàncies a l'espai
%  Exercicis: 74, 76, 86, 88
% ============================================================

\setmana{32}{Unitat 6 \textperiodcentered{} Angles i distàncies a l'espai}{25 d'abril de 2027}%
        {74, 76, 86 i 88}

\begin{errada}{Exercici 74}
  \ondiu{«c) Trobem la projecció ortogonal del punt sobre la recta.»}
  \hadedir{«b) Trobem la projecció ortogonal del punt sobre la recta.»}
  \motiu{L'exercici només té dos apartats; aquest és el b), el del punt
    $P(-1,2,13)$.}
\end{errada}

\begin{errada}{Exercici 86 b)}
  \ondiu{$\vec{n}=(2,-2,1)$}
  \hadedir{$\vec{n}=(2,-2,-1)$}
  \motiu{El coeficient de $z$ del pla $2x-2y-z+1=0$ és $-1$. El mòdul, $3$,
    i la distància, $1$ u, no canvien.}
\end{errada}

\begin{errada}{Exercici 88, equació}
  \ondiu{$|\lambda+4|=\sqrt{6}$}
  \hadedir{$\dfrac{|\lambda+4|}{\sqrt{6}}=\sqrt{6}$}
  \motiu{Hi falta el denominador $\sqrt{6}$ de la fórmula de la distància.
    La línia següent, $|\lambda+4|=6$, ja és correcta.}
\end{errada}

\begin{errada}{Exercici 88, segon punt}
  \ondiu{$\lambda+4=-6 \;\rightarrow\; \lambda=-8 \;\rightarrow\; P(-6,19,-8)$}
  \hadedir{$\lambda+4=-6 \;\rightarrow\; \lambda=-10 \;\rightarrow\; P(-8,23,-10)$}
  \motiu{De $\lambda+4=-6$ surt $\lambda=-10$. El punt $(-6,19,-8)$ és a
    distància $\tfrac{4}{\sqrt{6}}$ del pla, no a $\sqrt{6}$.}
\end{errada}

%\senseerrades{76}

%% ============================================================
%  Setmana 33 — Unitat 6 · Angles i distàncies a l'espai
%  Exercicis: 90, 92, 97, 115, 139
% ============================================================

\setmana{33}{Unitat 6 \textperiodcentered{} Angles i distàncies a l'espai}{2 de maig de 2027}%
        {90, 92, 97, 115 i 139}

\begin{errada}{Exercici 90 a)}
  \ondiu{$r:\ \dfrac{x}{1}=\dfrac{y+1}{-1}=\dfrac{z+1}{-1}$}
  \hadedir{$r:\ \dfrac{x}{1}=\dfrac{y+1}{-1}=\dfrac{z-1}{-1}$}
  \motiu{La recta passa per $A(0,-1,1)$, on $z=1$. La línia següent,
    $x=-z+1$, ja és correcta.}
\end{errada}

\begin{errada}{Exercici 90 b)}
  \ondiu{$d(A,\pi)=\dfrac{|3\cdot(-2)+5\cdot 5-2\cdot(-4)+7|}
    {\sqrt{3^2+5^2+(-2)^2}}=\dfrac{34}{\sqrt{38}}$ u}
  \hadedir{$d(A,\pi)=\dfrac{|3\cdot(-2)+5\cdot 1-2\cdot(-4)+7|}
    {\sqrt{3^2+5^2+(-2)^2}}=\dfrac{14}{\sqrt{38}}$ u}
  \motiu{La segona coordenada del punt $A(-2,1,-4)$ és $1$, no $5$. La
    distància és $\tfrac{14}{\sqrt{38}}=\tfrac{7\sqrt{38}}{19}\approx 2{,}27$ u,
    que és també la distància de $A$ al punt $Q$ que s'ha trobat.}
\end{errada}

\begin{errada}{Exercici 139 a)}
  \ondiu{$r:\ \dfrac{x-1}{1}=\dfrac{y-2}{2}=\dfrac{z-1}{3}$}
  \hadedir{$r:\ \dfrac{x-1}{1}=\dfrac{y-2}{1}=\dfrac{z-1}{3}$}
  \motiu{El vector director de la recta és el normal del pla,
    $\vec{n}=(1,1,3)$: el segon denominador és $1$.}
\end{errada}

\begin{errada}{Exercici 139 b)}
  \ondiu{$Q\left(\tfrac{1}{2},\,1,\,-\tfrac{1}{2}\right)$}
  \hadedir{$Q\left(\tfrac{5}{11},\,\tfrac{16}{11},\,-\tfrac{7}{11}\right)$}
  \motiu{Amb la recta correcta, $(1+t,\;2+t,\;1+3t)$, el pla dona $6+11t=0$,
    és a dir, $t=-\tfrac{6}{11}$.}
\end{errada}

\begin{errada}{Exercici 139 c)}
  \ondiu{$d(P,\pi)=\sqrt{\dfrac{7}{2}}$ u}
  \hadedir{$d(P,\pi)=\dfrac{6}{\sqrt{11}}=\dfrac{6\sqrt{11}}{11}$ u}
  \motiu{Amb la fórmula, $d=\dfrac{|1+2+3\cdot 1|}{\sqrt{1+1+9}}=
    \dfrac{6}{\sqrt{11}}\approx 1{,}81$ u. El valor de la solució és la
    distància de $P$ a un punt del pla que no és el peu de la perpendicular.}
\end{errada}

%\senseerrades{92, 97 i 115}

%% ============================================================
%  Setmana 34 — Unitat 11 · Integrals
%  Exercicis: 52, 54
% ============================================================

\setmana{34}{Unitat 11 \textperiodcentered{} Integrals}{9 de maig de 2027}%
        {52 i 54}

\begin{errada}{Exercici 52 g)}
  \ondiu{$\displaystyle\int (x-1)^3\,dx=\dfrac{(x-1)^3}{4}+k$}
  \hadedir{$\displaystyle\int (x-1)^3\,dx=\dfrac{(x-1)^4}{4}+k$}
  \motiu{En integrar una potència, l'exponent augmenta una unitat.}
\end{errada}

\begin{errada}{Exercici 54 a)}
  \ondiu{$\displaystyle\int\dfrac{4}{x-2}\,dx=4\ln|x+2|+k$}
  \hadedir{$\displaystyle\int\dfrac{4}{x-2}\,dx=4\ln|x-2|+k$}
  \motiu{El denominador és $x-2$.}
\end{errada}

\begin{errada}{Exercici 54 d)}
  \ondiu{$\displaystyle\int\dfrac{1}{(x+4)^2}\,dx=-\dfrac{1}{x-4}+k$}
  \hadedir{$\displaystyle\int\dfrac{1}{(x+4)^2}\,dx=-\dfrac{1}{x+4}+k$}
  \motiu{El denominador és $x+4$.}
\end{errada}

\begin{errada}{Exercici 54 f)}
  \ondiu{$\displaystyle\int\left(\dfrac{2}{(x+3)^3}+\dfrac{5}{(x-3)^2}\right)dx
    =\dfrac{-1}{(x+3)^2}-\dfrac{5}{x-3}+k$}
  \hadedir{$\displaystyle\int\left(\dfrac{2}{(x+3)^3}-\dfrac{5}{(x-3)^2}\right)dx
    =\dfrac{-1}{(x+3)^2}+\dfrac{5}{x-3}+k$}
  \motiu{A l'enunciat, el segon terme porta signe menys. La solució el copia
    amb signe més, i el resultat que dona és el d'aquesta altra integral.}
\end{errada}

%% ============================================================
%  Setmana 35 — Unitat 12 · La integral definida
%  Exercicis: 35, 39, 55, 79
% ============================================================

\setmana{35}{Unitat 12 \textperiodcentered{} La integral definida}{16 de maig de 2027}%
        {35, 39, 55 i 79}

\begin{errada}{Exercici 55 c)}
  \ondiu{$=\dfrac{9}{4}\sqrt[3]{9}+\dfrac{195}{12}$}
  \hadedir{$=\dfrac{27}{4}\sqrt[3]{9}+\dfrac{199}{12}$}
  \motiu{$\tfrac{3}{4}\sqrt[3]{9^4}=\tfrac{3}{4}\cdot 9\sqrt[3]{9}=
    \tfrac{27}{4}\sqrt[3]{9}$ \;i\; $18-\tfrac{3}{4}-\tfrac{2}{3}=\tfrac{199}{12}$.
    El valor és aproximadament $30{,}62$.}
\end{errada}

\begin{errada}{Exercici 55 g)}
  \ondiu{$\displaystyle\int_0^{\pi}\cos 2x\,dx$}
  \hadedir{$\displaystyle\int_0^{\pi}\cos^2 x\,dx$}
  \motiu{L'enunciat demana la integral de $\cos^2 x$, i la primitiva
    $\tfrac{x}{2}+\tfrac{\sin 2x}{4}$ i el resultat $\tfrac{\pi}{2}$ són els
    d'aquesta funció. La integral de $\cos 2x$ entre $0$ i $\pi$ valdria $0$.}
\end{errada}

%\senseerrades{35, 39 i 79}

%% ============================================================
%  Setmana 36 — Unitat 12 · La integral definida
%  Exercicis: 80, 88, 96, 111, 114
% ============================================================

\setmana{36}{Unitat 12 \textperiodcentered{} La integral definida}{23 de maig de 2027}%
        {80, 88, 96, 111 i 114}

\noindent
No hi ha errors.

%\senseerrades{80, 88, 96, 111 i 114}




\end{document}
